\documentclass[12pt]{article}
\usepackage{geometry}
\geometry{margin=1in}

% ----------------------------------------------------------
%  Linux Penguin Theme — Based on Recommended Colors
% ----------------------------------------------------------

% --- COLOR SETUP ---
\usepackage{xcolor}

\definecolor{cream}{RGB}{246,242,219}
\definecolor{teamred}{RGB}{209,57,64}
\definecolor{navy}{RGB}{40,46,87}
\definecolor{softwhite}{RGB}{250,250,230}
\definecolor{accentyellow}{RGB}{240,200,60}

% Extra utility colors
\definecolor{muted}{RGB}{120,125,140}
\definecolor{darkcream}{RGB}{230,225,205}

% ----------------------------------------------------------
%  FONT SETUP (Optional but Recommended)
%  Works best with XeLaTeX or LuaLaTeX
% ----------------------------------------------------------
\usepackage{fontspec}

% Main document font (Google fonts — upload to Overleaf)
% If you don't upload fonts, comment these two lines.
\setmainfont{SourceSans3}[
  Path=./,
  ItalicFont  = *-Italic,
  BoldFont    = *-Bold
]

\newfontfamily\headingfont{ScienceGothic}[
  Path=./,
]

% Font Awesome for icons
\usepackage{fontawesome5}

% ----------------------------------------------------------
%  SECTION STYLE
% ----------------------------------------------------------
\usepackage{titlesec}

% Section
\titleformat{\section}
  {\Large\headingfont\bfseries\color{navy}}
  {\thesection}
  {1em}
  {\rule{\linewidth}{1pt}\vspace{0.5em}\\}

% Subsection
\titleformat{\subsection}
  {\large\headingfont\bfseries\color{teamred}}
  {\thesubsection}
  {0.8em}
  {}

% Subsubsection
\titleformat{\subsubsection}
  {\bfseries\color{navy}}
  {\thesubsubsection}
  {0.6em}
  {}

% ----------------------------------------------------------
%  PAGE STYLE
% ----------------------------------------------------------
\usepackage{fancyhdr}
\pagestyle{fancy}

\fancyhf{} % clear all

\lhead{\textcolor{navy}{\headingfont \courseTitle}}
\rhead{\textcolor{teamred}{\thepage}}
\renewcommand{\headrule}{\color{teamred}\hrule height 1pt}

% ----------------------------------------------------------
%  HYPERLINK STYLE
% ----------------------------------------------------------
\usepackage{hyperref}

\hypersetup{
  colorlinks = true,
  linkcolor  = teamred,
  urlcolor   = navy,
  citecolor  = accentyellow
}

% ----------------------------------------------------------
%  CUSTOM BOX ENVIRONMENTS
% ----------------------------------------------------------
\usepackage{tcolorbox}
\tcbuselibrary{skins,breakable}

% Info box
\newtcolorbox{infobox}{
  enhanced,
  breakable,
  colback   = softwhite,
  colframe  = navy,
  coltext   = navy,
  boxrule   = 1pt,
  arc       = 4pt,
  left=6pt, right=6pt, top=6pt, bottom=6pt
}

% Warning box
\newtcolorbox{warningbox}{
  enhanced,
  breakable,
  colback   = accentyellow!30,
  colframe  = accentyellow!80!navy,
  coltext   = navy,
  boxrule   = 1pt,
  arc       = 4pt,
  left=6pt, right=6pt, top=6pt, bottom=6pt
}

% Highlight box
\newtcolorbox{highlightbox}{
  enhanced,
  breakable,
  colback   = teamred!10,
  colframe  = teamred,
  coltext   = navy,
  boxrule   = 1pt,
  arc       = 4pt,
  left=6pt, right=6pt, top=6pt, bottom=6pt
}

% ----------------------------------------------------------
%  ITEMIZE / ENUMERATE STYLE
% ----------------------------------------------------------
\usepackage{enumitem}

\setlist[itemize]{
  label=\textcolor{teamred}{\large\textbullet},
  leftmargin=1.2em,
  itemsep=4pt
}

\setlist[enumerate]{
  label=\textcolor{navy}{\arabic*.},
  leftmargin=1.4em,
  itemsep=4pt
}

% ----------------------------------------------------------
%  TITLE STYLE
% ----------------------------------------------------------
\usepackage{titling}

\pretitle{
  \begin{center}
  \headingfont\bfseries\LARGE\color{navy}
}
\posttitle{
  \end{center}
}

\preauthor{
  \begin{center}
  \color{teamred}
}
\postauthor{
  \end{center}
}

% ----------------------------------------------------------
%  OPTIONAL — SOFT PAGE BACKGROUND
%  Comment out if you want white pages
% ----------------------------------------------------------
\usepackage{pagecolor}
\pagecolor{cream}
\color{navy}

% ----------------------------------------------------------
% CUSTOM COMMANDS FOR TEMPLATE
% ----------------------------------------------------------
\newcommand{\courseTitle}{User Permissions and File System on Linux}
\newcommand{\courseDate}{17-12-2025}

% Icon commands using Font Awesome
\newcommand{\penguinicon}{\faLinux}
\newcommand{\importanticon}{\faExclamationCircle}
\newcommand{\treeicon}{\faTree}
\newcommand{\homeicon}{\faHome}
\newcommand{\configicon}{\faCogs}
\newcommand{\databaseicon}{\faDatabase}
\newcommand{\programicon}{\faCode}
\newcommand{\fileicon}{\faFile}
\newcommand{\hiddenicon}{\faEyeSlash}
\newcommand{\usericon}{\faUser}
\newcommand{\groupicon}{\faUsers}
\newcommand{\permissionicon}{\faLock}
\newcommand{\rooticon}{\faUserShield}
\newcommand{\linkicon}{\faLink}
\newcommand{\warningicon}{\faExclamationTriangle}
\newcommand{\coffeeicon}{\faCoffee}
\newcommand{\checkicon}{\faCheck}
\newcommand{\handicon}{\faHandPointRight}

% ----------------------------------------------------------
% END OF THEME
% ----------------------------------------------------------

\begin{document}

% Title with minimal space
\begin{center}
\headingfont\Huge\bfseries\color{navy}
\penguinicon\ \courseTitle\\[0.5em]
\Large\color{teamred}
\courseDate
\end{center}

\vspace{1em}
\begin{center}
\color{teamred}\rule{0.8\textwidth}{1pt}
\end{center}

\section*{\importanticon\ Why This Topic Matters (Seriously)}

Linux doesn't break because it's fragile.
Linux breaks because \textbf{permissions exist} and \textbf{you don't understand them yet}.

Permissions:

\begin{itemize}
    \item Protect your system
    \item Prevent malware from wrecking everything
    \item Make Linux powerful for servers, DevOps, security, and cloud
\end{itemize}

Once this clicks:

\begin{itemize}
    \item Errors make sense
    \item "Permission denied" stops ruining your day
    \item You stop nuking your system with \texttt{sudo rm -rf}
\end{itemize}

\section*{1 \treeicon\ Linux File System Hierarchy (The Big Picture)}

Linux doesn't use drive letters like Windows (\texttt{C:\textbackslash}, \texttt{D:\textbackslash}).
Instead, \textbf{everything starts at one place}:

\subsection*{\texttt{/} — The Root of Everything}

\texttt{/} is the \textbf{top of the file system tree}.

Think of it like:

\begin{infobox}
The trunk of a tree \treeicon\\
Everything grows from here.
\end{infobox}

\begin{highlightbox}
\texttt{\$ /}\\
\texttt{├── home}\\
\texttt{├── etc}\\
\texttt{├── var}\\
\texttt{├── usr}\\
\texttt{├── bin}\\
\texttt{├── sbin}\\
\texttt{└── tmp}
\end{highlightbox}

\begin{warningbox}
\warningicon\ You \textbf{do not} casually mess around in \texttt{/}.
\end{warningbox}

\subsection*{\texttt{/home} — Where Humans Live}

This is where \textbf{users live}.

\begin{highlightbox}
\texttt{\$ /home/chuck}\\
\texttt{\$ /home/student}
\end{highlightbox}

Inside your home directory:

\begin{itemize}
    \item Documents
    \item Downloads
    \item Config files
    \item SSH keys
    \item Hidden dot files
\end{itemize}

Your home is:

\begin{itemize}
    \item Owned by you
    \item Writable by you
    \item Mostly safe from system damage
\end{itemize}

\begin{infobox}
\handicon\ \textbf{Best Practice:}\\
If you're learning Linux, live in \texttt{/home}
\end{infobox}

\subsection*{\texttt{/etc} — Configuration Central}

This is one of the \textbf{most important directories}.

\texttt{/etc} contains:

\begin{itemize}
    \item System configuration files
    \item User settings
    \item Service configs
    \item Network configs
\end{itemize}

Examples:

\begin{highlightbox}
\texttt{\$ /etc/passwd}\\
\texttt{\$ /etc/shadow}\\
\texttt{\$ /etc/hosts}\\
\texttt{\$ /etc/sudoers}
\end{highlightbox}

\begin{warningbox}
\warningicon\ Editing files here can:
\begin{itemize}
    \item Break login
    \item Break networking
    \item Lock you out completely
\end{itemize}
\end{warningbox}

\begin{infobox}
\handicon\ \textbf{Rule:}\\
\texttt{/etc} = configuration, not programs
\end{infobox}

\subsection*{\texttt{/var} — Variable Data (Stuff That Changes)}

This is where Linux stores:

\begin{itemize}
    \item Logs
    \item Caches
    \item Spools
    \item Databases
\end{itemize}

Examples:

\begin{highlightbox}
\texttt{\$ /var/log}\\
\texttt{\$ /var/log/syslog}\\
\texttt{\$ /var/log/auth.log}
\end{highlightbox}

\begin{infobox}
\handicon\ If a system is acting weird:\\
\handicon\ \textbf{Try Checking \texttt{/var/log} or Any other related logs for your problem}
\end{infobox}

\subsection*{\texttt{/usr} — User Programs (Not Users)}

Confusing name.
This is \textbf{not user home directories}.

\texttt{/usr} contains:

\begin{itemize}
    \item Installed programs
    \item Libraries
    \item Shared resources
\end{itemize}

\begin{highlightbox}
\texttt{\$ /usr/bin}\\
\texttt{\$ /usr/lib}\\
\texttt{\$ /usr/share}
\end{highlightbox}

Most software you install ends up here.

\section*{2 \fileicon\ "Everything Is a File" (Linux Philosophy)}

This is where Linux becomes \textit{cool}.

In Linux:

\begin{itemize}
    \item Files
    \item Directories
    \item Devices
    \item Hardware
    \item Network interfaces
\end{itemize}

\handicon\ \textbf{Everything is a file}

Examples:

\begin{highlightbox}
\texttt{\$ /dev/sda      \# hard drive}\\
\texttt{\$ /dev/null     \# black hole}\\
\texttt{\$ /proc/cpuinfo \# CPU info}\\
\texttt{\$ /dev/tty      \# current terminal}\\
\texttt{\$ /dev/random   \# random number generator}\\
\texttt{\$ /dev/urandom  \# faster random number generator}\\
\texttt{\$ /proc/uptime       \# system uptime}\\
\texttt{\$ /proc/net/tcp      \# TCP connections}\\
\texttt{\$ /sys/class/net/eth0/address  \# network MAC address}\\
\texttt{\$ /sys/class/backlight/brightness  \# screen brightness}\\
\texttt{\$ /sys/class/power\_supply/BAT0/capacity  \# battery level}\\
\texttt{\$ /dev/loop0    \# loopback device (mount files as disks)}
\end{highlightbox}

Want hardware info?

\begin{highlightbox}
\texttt{\$ cat /proc/cpuinfo}
\end{highlightbox}

Want to silence output?

\begin{highlightbox}
\texttt{\$ command > /dev/null}
\end{highlightbox}

\begin{infobox}
\handicon\ This design is WHY:
\begin{itemize}
    \item Linux is scriptable
    \item Linux is automatable
    \item Linux runs the internet
\end{itemize}
\end{infobox}

\section*{3 \hiddenicon\ Dot Files (Hidden Power)}

Files that start with a dot \texttt{.} are \textbf{hidden}.

Example:

\begin{highlightbox}
\texttt{\$ .bashrc}\\
\texttt{\$ .profile}\\
\texttt{\$ .gitconfig}\\
\texttt{\$ .ssh/}
\end{highlightbox}

They usually store:

\begin{itemize}
    \item Shell configs
    \item App preferences
    \item User settings
\end{itemize}

List them:

\begin{highlightbox}
\texttt{\$ ls -a}
\end{highlightbox}

Why hide them?

\begin{itemize}
    \item Prevent clutter
    \item Prevent accidental deletion
\end{itemize}

\begin{infobox}
\handicon\ \textbf{Important:}\\
Deleting dot files can reset your environment or break tools.
\end{infobox}

\section*{4 \usericon\ File Ownership (Who Owns What?)}

Every file has:

\begin{itemize}
    \item \textbf{Owner}
    \item \textbf{Group}
    \item \textbf{Permissions}
\end{itemize}

Check with:

\begin{highlightbox}
\texttt{\$ ls -l}
\end{highlightbox}

Example:

\begin{highlightbox}
\texttt{\$ -rwxr-xr-- 1 ahmad staff script.sh}
\end{highlightbox}

Breakdown:

\begin{infobox}
\texttt{Owner: ahmad}\\
\texttt{Group: staff}
\end{infobox}

Ownership matters because:

\begin{itemize}
    \item Owners control permissions
    \item Groups allow shared access
    \item Others are everyone else
\end{itemize}

\subsection*{Changing Ownership}

\begin{highlightbox}
\texttt{\$ chown user file}\\
\texttt{\$ chown user:group file}
\end{highlightbox}

Example:

\begin{highlightbox}
\texttt{\$ sudo chown root:root config.conf}
\end{highlightbox}

\begin{warningbox}
\warningicon\ Wrong ownership = broken apps
\end{warningbox}

\section*{5 \permissionicon\ Permissions (rwx) — The Core Skill}

Permissions are \textbf{everything}.

Three types:

\begin{itemize}
    \item \textbf{r} = read
    \item \textbf{w} = write
    \item \textbf{x} = execute
\end{itemize}

Three levels:

\begin{itemize}
    \item Owner
    \item Group
    \item Others
\end{itemize}

Example:

\begin{highlightbox}
\texttt{\$ -rwxr-xr--}
\end{highlightbox}

Break it down:

\begin{infobox}
\texttt{Owner:  rwx}\\
\texttt{Group:  r-x}\\
\texttt{Others: r--}
\end{infobox}

\subsection*{Permission Numbers (Octal Mode)}

This is where beginners panic—but don't.

Values:

\begin{infobox}
\texttt{r = 4}\\
\texttt{w = 2}\\
\texttt{x = 1}
\end{infobox}

Add them:

\begin{infobox}
\texttt{7 = rwx}\\
\texttt{6 = rw-}\\
\texttt{5 = r-x}\\
\texttt{4 = r--}
\end{infobox}

Example:

\begin{highlightbox}
\texttt{\$ chmod 755 script.sh}
\end{highlightbox}

Meaning:

\begin{infobox}
\texttt{Owner: 7 (rwx)}\\
\texttt{Group: 5 (r-x)}\\
\texttt{Others: 5 (r-x)}
\end{infobox}

\begin{infobox}
\handicon\ You will see \texttt{755} and \texttt{644} everywhere.
\end{infobox}

\subsection*{Making a Script Executable}

\begin{highlightbox}
\texttt{\$ chmod +x script.sh}
\end{highlightbox}

Without execute permission:

\begin{highlightbox}
\texttt{\$ ./script.sh}\\
\texttt{\# Permission denied}
\end{highlightbox}

\section*{6 \rooticon\ Root User (Unlimited Power)}

Root is:

\begin{itemize}
    \item The system administrator
    \item The power mode user
    \item The "I can change everything" account
\end{itemize}

Root can:

\begin{itemize}
    \item Read any file
    \item Delete any file
    \item Change any permission
\end{itemize}

\subsection*{\texttt{sudo} — Temporary Root}

\begin{highlightbox}
\texttt{\$ sudo command}
\end{highlightbox}

This:

\begin{itemize}
    \item Executes command as root
    \item Logs actions
    \item Requires password
\end{itemize}

\begin{infobox}
\handicon\ \textbf{Best practice:}\\
Use \texttt{sudo}, not permanent root login
\end{infobox}

\subsection*{\texttt{su} — Switch User}

\begin{highlightbox}
\texttt{\$ su}\\
\texttt{\$ su username}
\end{highlightbox}

This:

\begin{itemize}
    \item Switches shell to another user
    \item Requires target user password
\end{itemize}

\begin{warningbox}
\warningicon\ Dangerous if misused
\end{warningbox}

\subsection*{\texttt{/etc/sudoers} — Who Can Use sudo}

DO NOT edit directly.

Always use:

\begin{highlightbox}
\texttt{\$ sudo visudo}
\end{highlightbox}

Example entry:

\begin{highlightbox}
\texttt{ahmad ALL=(ALL) ALL}
\end{highlightbox}

\begin{warningbox}
\warningicon\ Wrong syntax here => no sudo => bad day.
\end{warningbox}

\section*{7 \linkicon\ Hard Links vs Soft Links (Symbolic Links)}

Links are \textbf{references} to files.

\subsection*{Hard Links}

\begin{itemize}
    \item Point to the same inode
    \item File exists in multiple locations
    \item Deleting one doesn't remove data
\end{itemize}

\begin{highlightbox}
\texttt{\$ ln file1 file2}
\end{highlightbox}

Limitations:

\begin{itemize}
    \item Same filesystem only
    \item Can't link directories (usually)
\end{itemize}

\subsection*{Soft Links (Symlinks)}

\begin{itemize}
    \item Pointer to file path
    \item Like Windows shortcuts
\end{itemize}

\begin{highlightbox}
\texttt{\$ ln -s target linkname}
\end{highlightbox}

Example:

\begin{highlightbox}
\texttt{\$ ln -s /var/log/syslog syslog\_link}
\end{highlightbox}

If target is deleted:
\begin{warningbox}
\warningicon\ Link breaks
\end{warningbox}

\begin{infobox}
\handicon\ Most commonly used type.
\end{infobox}

\section*{8 \warningicon\ Common Beginner Mistakes (Read This Twice)}

This section will save your system.

\subsection*{\warningicon\ Running Everything with sudo}

Bad:

\begin{highlightbox}
\texttt{\$ sudo npm install}\\
\texttt{\$ sudo python script.py}
\end{highlightbox}

Why bad?

\begin{itemize}
    \item Files owned by root
    \item Permission nightmares
    \item Security risk
\end{itemize}

\begin{infobox}
\handicon\ Use sudo \textbf{only when required}
\end{infobox}

\subsection*{\warningicon\ \texttt{chmod 777} Everything}

\begin{highlightbox}
\texttt{\$ chmod 777 file}
\end{highlightbox}

This means:

\begin{infobox}
Everyone can do everything
\end{infobox}

Bad for:

\begin{itemize}
    \item Security
    \item Best practices
    \item Your future job interviews \faSmile
\end{itemize}

\subsection*{\warningicon\ Editing System Files Without Backup}

Before:

\begin{highlightbox}
\texttt{\$ sudo cp config.conf config.conf.bak}
\end{highlightbox}

Always.

\subsection*{\warningicon\ Deleting Random Stuff in \texttt{/etc} or \texttt{/usr}}

If you don't know what it does:
\handicon\ \textbf{Don't delete it}

\subsection*{\warningicon\ Not Reading Error Messages}

Linux tells you:

\begin{itemize}
    \item What failed
    \item Why it failed
    \item What permission is missing
\end{itemize}

Read it.

\section*{9 \handicon\ Real-World Permission Scenarios}

\begin{itemize}
    \item \textbf{"Permission denied" when running script}
    \begin{itemize}
        \item Possible Missing execute bit
    \end{itemize}
    
    \item \textbf{App won't start}
    \begin{itemize}
        \item Maybe Wrong ownership in config directory
    \end{itemize}
    
    \item \textbf{Can't write to directory}
    \begin{itemize}
        \item Check group permissions
    \end{itemize}
    
    \item \textbf{Service won't start}
    \begin{itemize}
        \item Maybe Config owned by wrong user
    \end{itemize}
\end{itemize}

\section*{10 \penguinicon\ Final Linux Mindset Shift}

Linux is not hard.

Linux is:

\begin{itemize}
    \item Honest
    \item Strict
    \item Predictable
\end{itemize}

Permissions are not obstacles.
They are \textbf{guardrails}.

Once you understand:

\begin{itemize}
    \item File hierarchy
    \item Ownership
    \item Permissions
    \item Root access
\end{itemize}

You unlock:

\begin{itemize}
    \item Servers
    \item DevOps
    \item Security
    \item Cloud
    \item Automation
\end{itemize}

\subsection*{What You Should Practice Next}

\begin{itemize}
    \item \checkicon\ \texttt{ls -l} everywhere
    \item \checkicon\ \texttt{chmod} with numbers
    \item \checkicon\ \texttt{chown} safely
    \item \checkicon\ Explore \texttt{/etc} (read-only)
    \item \checkicon\ Break things \textbf{inside a VM}
\end{itemize}

\begin{infobox}
If you want, next we can:
\begin{itemize}
    \item Do a \textbf{hands-on lab}
    \item Build a \textbf{permission challenge}
    \item Simulate \textbf{real server mistakes}
    \item Or connect this to \textbf{Docker \& servers}
\end{itemize}
\end{infobox}

Now go grab that coffee \coffeeicon\\
You just leveled up in Linux. \penguinicon\ \faFire

% Footer
\vspace{\fill}
\begin{center}
\color{muted}
All rights reserved for \texttt{TogetherTeam} [NOT For Commercial Use] !!\\
contact the author via email \href{mailto:khaledmhfz2004@gmail.com}{\texttt{khaledmhfz2004@gmail.com}}
\end{center}
\end{document}